\section{INTRODUCTION}

Selenium is an essential trace element for mammals that participates in the biosynthesis of selenoproteins such as glutathione peroxidase, playing a pivotal role in maintaining organismal redox homeostasis. Compared to inorganic selenium species such as selenite and selenate, zero-valent selenium nanoparticles (SeNPs) exhibit higher bioavailability and lower acute toxicity. Benefiting from their favorable biocompatibility and diverse biological activities, SeNPs hold broad application prospects across biomedical fields, including anti-tumor, antibacterial, antidiabetic, and antioxidant therapeutics \cite{hosnedlova2018nanoselenium,bisht2022senps,sowmya2024nanoselenium,karthik2024senps}. However, selenium possesses a narrow therapeutic index with a limited margin between effective and toxic doses; even within pharmacologically active dose regimens, SeNPs can still induce adverse outcomes such as oxidative damage and cytotoxicity \cite{wang2020prooxidation}. Furthermore, excessive exposure may precipitate systemic toxicities, including fever, nausea, respiratory, cardiovascular, and gastrointestinal impairments, and even renal failure \cite{hosnedlova2018nanoselenium}, highlighting the necessity and urgency of comprehensive biosafety evaluations.

The biotoxicity of nanomaterials is collectively governed by multifaceted factors such as particle size, morphology, surface properties, and chemical composition. Marked non-linear coupling characteristics exist among these physicochemical parameters, compounding the complexity of toxicity prediction and mechanistic elucidation \cite{ji2025rosml}. Conventional nanotoxicity assessments predominantly rely on in vitro cell-based assays and in vivo animal models (e.g., rat, mouse, and zebrafish models); nevertheless, such experimental paradigms generally suffer from protracted testing cycles, prohibitive costs, and constrained throughput \cite{ahmadi2024toxicity}. With the advancement of computational toxicology, nano-quantitative structure--activity/toxicity relationship (nano-QSAR/QSTR) methodologies have progressively emerged as indispensable complementary tools in nanotoxicity research \cite{roy2023tio2}. In recent years, machine learning algorithms such as Random Forest (RF), Support Vector Machines (SVM), and Neural Networks (NN) have been widely deployed to predict the cytotoxicity of metal and metal oxide nanoparticles \cite{zhang2023qsar}, frequently coupled with clustering analysis and feature selection techniques to further enhance model stability, predictive accuracy, and generalizability \cite{sang2022mlcytotoxicity}.

Dedicated investigations on selenium nanoparticles have revealed that their cytotoxic effects are intimately linked to factors including particle size, exposure concentration, and surface functionalization. Regarding particle size, as SeNP dimensions decrease and specific surface area increases, cellular uptake efficiency and biological activity are typically amplified, accompanied by an elevated risk of potential toxicity \cite{zhang2023senpsbio}. In terms of concentration, Wang et al. \cite{wang2020prooxidation} demonstrated that exposing rat hepatocytes to 24~$\mu$M and 48~$\mu$M SeNPs for 24~h reduced cell viability to 87.76\% and 51.37\%, respectively, manifesting an evident dose-dependent cytotoxic response. Surface modification likewise exerts a substantial impact on the biosafety and toxicological behavior of SeNPs. For instance, surface functionalization with polysaccharides such as Spirulina polysaccharide (SPS), Gracilaria lemaneiformis polysaccharide (GLP), and Astragalus polysaccharide (AASP) can potentiate the selective killing efficacy of SeNPs against tumor cells while mitigating their toxicity toward normal cells \cite{zhang2023senpsbio,ji2023astragalus}; conversely, quaternized chitosan--gallic acid-modified SeNPs exhibit pronounced cytotoxicity in Colo357 and HaCaT cells \cite{lunkov2023chitosan}. These findings substantiate that the toxicological profile of SeNPs exhibits pronounced physicochemical dependency that defies simplistic empirical extrapolation, underscoring the necessity of establishing toxicity prediction models tailored to specific physicochemical descriptors of selenium nanoparticles.

Recently, machine learning and deep learning have exhibited salient strengths in high-dimensional feature mining, complex non-linear relationship modeling, and pattern recognition, offering innovative technological routes for computational toxicology \cite{hosni2025mlnp}. Shirokii et al. utilized correlation-based feature selection and grid search hyperparameter optimization to elevate the 10-fold cross-validation $Q^2$ of a quantitative cytotoxicity prediction model for inorganic nanomaterials to 0.86 \cite{shirokii2023quantitative}. Furthermore, ensemble learning architectures such as LightGBM have displayed superior predictive performance over conventional machine learning models in distinguishing cytotoxic from non-cytotoxic nanoparticles \cite{hosni2025mlnp}, while hybrid modeling frameworks such as Boruta-RF-ARM have further pinpointed critical toxicological determinants, including exposure concentration, core--shell composition, and cell lineage \cite{yildirim2025rheumatoid}. Despite these advancements in nanotoxicity modeling, prevailing studies have largely focused on organic small molecules or generalized nanomaterial categories. Dedicated investigations into the cytotoxic profile of selenium nanoparticles remain scarce, and domain-specific predictive models alongside feature engineering frameworks for SeNPs are still lacking. To bridge this gap, this study establishes an analytical framework for selenium nanoparticle cytotoxicity prediction comprising multi-model benchmarking, enhanced feature engineering, and automated feature selection. Taking essential variables—such as concentration, particle size, surface modification, and material chemical properties—as inputs, this work systematically evaluates the applicability and predictive efficacy of diverse models across both generalized nanoparticle benchmarks and dedicated selenium nanoparticle datasets, aiming to provide a methodological foundation for SeNP biosafety assessment, experimental candidate prioritization, and the future development of intelligent prediction platforms.
